% Appendix -- encoder-track details. Owner: section agent B.
% Sources: clarity/encoder.py (Cfg defaults, optimiser, schedule), clarity/decide.py (tau grid),
% clarity/experiments/*/common.args and lane files, clarity/reports/02_experiment_log.md (E0-E12),
% mideval/RESULTS.md and mideval/analysis/mideval_analysis.txt, clarity/reports/raw/E11_replication.txt.
% Cells start with \raggedright and rows end with \tabularnewline (no array package needed).
\section{Encoder-track details}
\label{sec:app-encoder}

\paragraph{Hyperparameters.}
Table~\ref{tab:app-enc-hparams} lists the settings, taken from
\texttt{clarity/encoder.py}, \texttt{clarity/decide.py} and the experiment argument
files. Baseline and final models differ only in the input and the number of epochs.
One implementation detail mattered: \texttt{transformers} 5.x loads a checkpoint in
its stored dtype, and DeBERTa-v3's published weights are fp16, in which its
attention overflows. With gradient clipping masking the overflow, the model trained
smoothly to the label prior (loss 1.887, the label entropy) and predicted Explicit
for every item. We therefore load the weights in fp32 explicitly.

\begin{table}[t]
\centering
\footnotesize
\begin{tabular}{@{}lp{0.64\columnwidth}@{}}
\toprule
Setting & Value \tabularnewline
\midrule
Backbone & \raggedright \texttt{microsoft/deberta-v3-large}, all weights trained \tabularnewline
Head & \raggedright \texttt{[CLS]} $\to$ dense + tanh $\to$ dropout 0.1 $\to$ linear (9) \tabularnewline
Input, baseline & \raggedright sub-question (${\le}96$ tokens) [SEP] answer; 512 tokens \tabularnewline
Input, final & \raggedright ``Sub-question: \ldots{} Full question: \ldots'' (${\le}256$ tokens) [SEP] answer; 1,024 tokens \tabularnewline
Truncation & \raggedright answer cut from the right \tabularnewline
Loss & \raggedright cross-entropy, no smoothing, no class weights \tabularnewline
Optimiser & \raggedright AdamW, $\beta=(0.9, 0.98)$, $\epsilon=10^{-6}$, weight decay 0.01 (none on biases and LayerNorm) \tabularnewline
Learning rate & \raggedright $10^{-5}\times0.95^{k}$ for a layer $k$ levels below the head; head and pooler $10^{-4}$ \tabularnewline
Schedule & \raggedright cosine, warm-up 10\% of steps \tabularnewline
Batch & \raggedright 16, no accumulation; gradient-norm clipping 1.0 \tabularnewline
Epochs & \raggedright 8 (baseline) or 16 \tabularnewline
Precision & \raggedright fp32 weights, bf16 autocast; gradient checkpointing \tabularnewline
Epoch selection & \raggedright best macro-F1 on the 345-row train slice; last epoch for all-of-train runs \tabularnewline
Seeds & \raggedright 0--9 (all of train: 0--4) \tabularnewline
Logit adjustment & \raggedright $\tau\in[-1,2]$, 61-point grid; prior = train label frequencies; 5-fold nested CV on dev \tabularnewline
Hardware & \raggedright one shared RTX PRO 6000 (96\,GB); peak memory about 9\,GB at 512 tokens, 15--17\,GB at 1,024 \tabularnewline
\bottomrule
\end{tabular}
\caption{DeBERTa-v3-large training and decision settings.}
\label{tab:app-enc-hparams}
\end{table}

\paragraph{Negative results.}
Table~\ref{tab:app-negative} lists each idea that did not help, with its control and
the reason the experiment log gives. Rows marked ``5 seeds'' compare against the
baseline's or the full-question model's seeds 0--4. The baseline's 5-seed system
(0.438) is the favourable draw discussed in \S\ref{sec:encoder}.

\begin{table*}[t]
\centering
\footnotesize
\begin{tabular}{@{}p{0.25\textwidth}p{0.17\textwidth}p{0.14\textwidth}p{0.35\textwidth}@{}}
\toprule
Idea & Result (dev \Stwo) & Control & Why it did not help \tabularnewline
\midrule
\raggedright Second DeBERTa re-ranks the baseline's top-3 labels from their definitions (5 seeds) &
\raggedright 0.354\,$\pm$\,0.003 with the baseline prior; 0.329 alone &
\raggedright 0.365 (baseline ensemble, argmax) &
\raggedright Training loss moved only 1.06$\to$0.88 in 8 epochs (1.10 = chance among three); learning
what a label means from its definition is harder than classifying. It fixed 25
baseline errors and broke 37 correct answers. \tabularnewline
\raggedright Full question + Balanced Softmax + focal loss, 8 epochs (5 seeds) &
\raggedright system 0.379; single \pmsd{0.322}{0.021} &
\raggedright 0.438; \pmsd{0.337}{0.044} &
\raggedright Over-corrected towards rare classes (Explicit 0.679$\to$0.319, Dodging
0.475$\to$0.081) and used up the gain of logit adjustment; most of the deficit was
the undertrained input. \tabularnewline
\raggedright Post-hoc routing: clarity level first, then leaf &
\raggedright 0.369 &
\raggedright 0.365 (ensemble, argmax) &
\raggedright The model's own top 3 contains an acceptable label for 89.6\% of items,
against 81.2\% for the taxonomy branch (3.60 labels). \tabularnewline
\raggedright Factorised head $p(\text{level})\,p(\text{leaf}\mid\text{level})$ (5 seeds) &
\raggedright system 0.370; single \pmsd{0.320}{0.038} &
\raggedright 0.438; \pmsd{0.337}{0.044} &
\raggedright \Sone{} steadier across seeds (sd 0.010 against 0.025 for the level
read-out) but not better; \Stwo{} lower. \tabularnewline
\raggedright Non-Reply gate + branch specialist encoders &
\raggedright 3-seed system 0.390; flat model's gate + specialists $-$0.009 per model
with the rule (5 seeds, 2/5) &
\raggedright 0.418 (flat, 3 seeds) &
\raggedright The gate memorises the binary split (training loss $\approx$0.03), so soft
routing becomes hard. The specialists passed the 3-seed screen and were level on 5
seeds. \tabularnewline
\raggedright Boundary experts for Implicit/Dodging, Explicit/Implicit and General/Deflection (3 seeds) &
\raggedright +0.002 per model (1/3); \Sone{} $-$0.004 &
\raggedright flat model, same seeds &
\raggedright They seldom overturn a decision the flat model makes between the same two
labels from the same input. \tabularnewline
\raggedright Uniform model soup of the 10 full-question models &
\raggedright 0.267 (greedy soup kept 1 model: 0.381) &
\raggedright 0.384 (single model) &
\raggedright Seeds differ in the random head initialisation and the data order, so
their weights do not lie in one basin; averaging any two already lowered the slice
score. \tabularnewline
\raggedright Nine per-class multipliers, nested CV (5-seed baseline ensemble) &
\raggedright 0.394 (in-fold minus held-out gap +0.098) &
\raggedright 0.438 (logit adjustment) &
\raggedright Nine parameters overfit 308 items; below argmax on two of the three
2-annotator reference sets. \tabularnewline
\raggedright Temperature + nine biases, fitted on one dev half and scored on the other &
\raggedright 0.293--0.461 over 4 configurations $\times$ 2 directions &
\raggedright logit adjustment &
\raggedright Erratic: the same rule swings by up to 0.148 between the two halves.
\tabularnewline
\raggedright Mixed-input ensemble: full question and sub-question + answer, both 16 epochs, seeds 5--9 &
\raggedright 0.403 &
\raggedright 0.405 (10 full-question models) &
\raggedright The idea came from one set of seeds (0.498) and did not replicate on fresh
seeds. \tabularnewline
\bottomrule
\end{tabular}
\caption{Encoder-track ideas that did not help, all on dev \Stwo. ``System'' is the
seed ensemble with logit adjustment ($\tau$ by nested CV); ``per model'' is a mean of
per-seed paired differences. Sources: \texttt{clarity/reports/02\_experiment\_log.md}
(E4--E12) and \texttt{mideval/RESULTS.md}.}
\label{tab:app-negative}
\end{table*}
